% =====================================================================
% APPENDIX (revised 2026-11-06: fork labels, 32B-original harness note, 1.x settings check, external tools,
% prior-work section for the new arXiv paper)
% Revised 2026-11-06 (pm): all 42 cells judged; "Pending cells" section removed; new sections app:reexport
% (selection-cap re-export provenance and validation) and app:external (external-tool provenance).
% Revised 2026-11-08 (ARS review): new app:seednoise (R1); added-cell marker double dagger (S3); app:reexport notes the
% 2.0.0 MLP range (R12); app:versions states the division of claims with the companion study (R2).
% =====================================================================
\section{Capability per Benchmark}
\label{app:capability}

% src (tests_capability.csv family original_design: 23 cells x 5 benchmarks, 0 Holm-significant; family added:
%      12 cells x 5, 1 Holm-significant (qwen2.5-7b__obliteratus-spectral_cascade mmlu_pro -7.0, p_holm 0.0003))
\begin{table}[!htbp]
\centering\footnotesize
\caption{Capability per benchmark. Rows ``Original'' give the original model's scores (\%); other rows
give the paired per-item change from that original (pp) and the change of the mean of the five benchmarks with a
95\% interval. None of the 115 per-benchmark changes of the main design is significant after Holm correction (exact
McNemar); of the 60 changes of the cells added after it ($^{\ddagger}$; separate Holm family), one is (bold: OBLITERATUS
\texttt{spectral\_cascade}, MMLU-Pro). $^{\dagger}$Unpaired comparison: the Qwen2.5-32B original was evaluated with an earlier harness revision
whose MBPP per-item samples predate the patched answer extractor; only its MBPP summary score was rescored from the
logged outputs, so its per-item scores are not available (\S\ref{sec:procedure}). SOM and J-space are plug-ins of
the companion study's project fork, not part of Heretic~2.0. One item is worth 0.1~pp (MMLU-Pro) to 0.61~pp
(HumanEval). The added cells are discussed in \S\ref{sec:res-kl} and \S\ref{sec:res-external}.}
\label{tab:capability}
\setlength{\tabcolsep}{4pt}
\resizebox{\textwidth}{!}{\begin{tabular}{@{}llrrrrrr@{}}
\toprule
Model & Condition & MMLU-Pro & GSM8K & IFEval & HumanEval & MBPP & Mean of five [95\% CI] \\
\midrule
Qwen2.5-7B & Original (score, \%) & 59.5 & 85.8 & 71.2 & 84.8 & 63.2 & 72.9 \\
 & 1.x general & +0.4 & $-$2.0 & $-$0.2 & $-$0.6 & +0.6 & $-$0.36 [$-$1.35, +0.63] \\
 & 1.x code & $-$1.4 & +3.2 & 0.0 & +0.6 & +1.6 & +0.80 [$-$0.41, +2.01] \\
 & 2.0 DiM & $-$1.4 & +0.2 & +0.6 & +1.2 & +0.6 & +0.23 [$-$0.85, +1.32] \\
 & 2.0 MPOA & $-$1.6 & $-$0.8 & $-$0.6 & +1.8 & 0.0 & $-$0.23 [$-$1.15, +0.70] \\
 & SOM (fork plug-in) & +0.3 & $-$1.8 & $-$0.6 & $-$1.2 & +1.6 & $-$0.33 [$-$1.44, +0.77] \\
 & J-space (fork plug-in) & +1.0 & +0.6 & +1.5 & 0.0 & +1.2 & +0.86 [$-$0.11, +1.82] \\
 & 2.0 ARA & $-$1.7 & +1.0 & $-$0.9 & +2.4 & +1.6 & +0.48 [$-$0.76, +1.72] \\
 & OBLITERATUS basic$^{\S}$ & $-$0.6 & +0.4 & $-$0.9 & $-$2.4 & $-$0.4 & $-$0.79 [$-$1.91, +0.33] \\
 & OBLITERATUS advanced$^{\S}$ & $-$3.4 & +1.2 & $-$2.2 & $-$1.8 & +1.4 & $-$0.97 [$-$2.30, +0.36] \\
 & OBLITERATUS spectral cascade$^{\S}$ & \textbf{$-$7.0} & $-$4.4 & $-$5.2 & $-$9.1 & $-$0.2 & $-$5.18 [$-$6.90, $-$3.47] \\
 & jwest33 recommended$^{\S}$ & $-$1.1 & $-$0.8 & +0.6 & $-$1.2 & +1.8 & $-$0.15 [$-$1.38, +1.08] \\
\addlinespace[2pt]
Qwen2.5-Coder-7B & Original (score, \%) & 41.6 & 83.4 & 58.8 & 89.6 & 72.2 & 69.1 \\
 & 1.x general & +0.5 & $-$1.4 & +0.6 & $-$0.6 & $-$0.8 & $-$0.35 [$-$1.18, +0.48] \\
 & 1.x code & +0.3 & $-$1.2 & +1.1 & +0.6 & +0.4 & +0.24 [$-$0.72, +1.21] \\
\addlinespace[2pt]
Qwen2.5-14B & Original (score, \%) & 67.4 & 84.2 & 79.9 & 82.3 & 71.0 & 77.0 \\
 & 1.x general & $-$1.7 & +1.2 & +1.7 & $-$1.2 & $-$0.6 & $-$0.13 [$-$1.04, +0.78] \\
 & 1.x code & $-$1.0 & +2.2 & $-$1.1 & $-$1.2 & +0.4 & $-$0.15 [$-$1.05, +0.76] \\
\addlinespace[2pt]
Qwen2.5-Coder-14B & Original (score, \%) & 49.7 & 88.8 & 65.8 & 90.9 & 74.8 & 74.0 \\
 & 1.x general & +0.4 & +0.2 & 0.0 & 0.0 & +0.6 & +0.24 [$-$0.59, +1.07] \\
 & 1.x code & +0.8 & $-$1.2 & +1.1 & +1.2 & +0.4 & +0.47 [$-$0.38, +1.31] \\
\addlinespace[2pt]
Qwen2.5-32B & Original (score, \%) & 71.1 & 85.4 & 78.9 & 89.6 & 77.6 & 80.5 \\
 & 1.x general & +0.4 & +0.4 & +0.4 & $-$0.6 & +0.4$^{\dagger}$ & +0.19 [$-$1.06, +1.45] \\
 & 1.x general, KL $\le$ 0.1$^{\S}$ & +0.1 & +0.8 & 0.0 & $-$0.6 & $-$0.2$^{\dagger}$ & +0.02 [$-$1.27, +1.31] \\
 & 1.x code & $-$0.1 & +0.2 & +1.3 & $-$2.4 & $-$0.2$^{\dagger}$ & $-$0.25 [$-$1.61, +1.11] \\
 & 2.0 DiM & $-$1.2 & +0.2 & +0.7 & +0.6 & $-$0.6$^{\dagger}$ & $-$0.05 [$-$1.35, +1.25] \\
 & 2.0 MPOA & $-$0.6 & $-$0.4 & +1.8 & 0.0 & $-$1.0$^{\dagger}$ & $-$0.03 [$-$1.30, +1.24] \\
 & SOM (fork plug-in) & 0.0 & +0.2 & +0.6 & 0.0 & +0.8$^{\dagger}$ & +0.31 [$-$0.92, +1.54] \\
 & 2.0 ARA & $-$0.8 & +1.2 & +0.9 & $-$4.9 & $-$0.6$^{\dagger}$ & $-$0.83 [$-$2.30, +0.64] \\
 & huihui-ai & $-$2.0 & +1.4 & +0.6 & $-$2.4 & $-$2.2$^{\dagger}$ & $-$0.94 [$-$2.49, +0.62] \\
 & zetasepic v2 & $-$1.8 & +1.2 & +2.2 & $-$3.0 & $-$0.8$^{\dagger}$ & $-$0.45 [$-$1.89, +1.00] \\
\addlinespace[2pt]
Qwen2.5-Coder-32B & Original (score, \%) & 62.7 & 92.0 & 80.0 & 90.9 & 80.0 & 81.1 \\
 & 1.x general & +1.3 & $-$0.4 & +0.4 & +1.2 & 0.0 & +0.50 [$-$0.21, +1.20] \\
 & 1.x code & +1.2 & +0.2 & $-$0.4 & +0.6 & $-$0.2 & +0.29 [$-$0.34, +0.92] \\
\addlinespace[2pt]
Qwen3.5-4B$^{\S}$ & Original (score, \%) & 70.6 & 91.8 & 82.3 & 79.9 & 62.6 & 77.4 \\
 & 2.0 DiM & $-$0.5 & $-$0.6 & +0.2 & +1.8 & $-$0.8 & +0.02 [$-$1.16, +1.21] \\
 & 2.0 MPOA & $-$1.4 & $-$0.2 & 0.0 & 0.0 & $-$1.0 & $-$0.52 [$-$1.51, +0.47] \\
 & SOM (fork plug-in) & $-$1.9 & +0.8 & $-$0.6 & 0.0 & $-$1.4 & $-$0.61 [$-$1.77, +0.55] \\
 & J-space (fork plug-in) & $-$0.2 & $-$0.2 & $-$0.9 & +3.0 & $-$1.4 & +0.06 [$-$1.16, +1.29] \\
 & 2.0 ARA & $-$1.2 & $-$1.6 & $-$3.0 & $-$0.6 & $-$1.4 & $-$1.55 [$-$3.11, 0.00] \\
 & jwest33, as released$^{\S}$ & $-$1.9 & +0.2 & $-$0.4 & +1.2 & $-$1.0 & $-$0.37 [$-$1.73, +0.99] \\
 & jwest33 + hybrid fix$^{\S}$ & $-$0.4 & $-$0.6 & $-$3.1 & +3.7 & $-$0.6 & $-$0.22 [$-$1.52, +1.08] \\
\bottomrule
\multicolumn{8}{@{}l}{\footnotesize Bold: per-benchmark change significant after Holm correction (pre-stated design: 115 tests; cells added$^{\S}$ after it: separate family of 60 tests).} \\
\multicolumn{8}{@{}l}{\footnotesize $^{\dagger}$ Unpaired test on summary scores (per-item scores of the original not reproducible).} \\
\end{tabular}
}
\end{table}

\section{Data Checks and Provenance}
\label{app:checks}

% src (RESULTS.md "Coverage" and "Anomalies and checks"; cells.csv adapter_sha256_ok, tensors_*; assemble_p30_log.txt;
%      reexport_1x/validation_7b/SUMMARY.json "original_merged_equals_base_plus_original_adapter")
\texttt{assemble\_p30.py} verifies each cell before analysis and lists every deviation in
\texttt{analysis/RESULTS.md}. For the analysis reported here:
\begin{itemize}
\item All 42 judged cells: 1{,}322 responses and judgments with matching identifiers and item hash; all judges
present on every row; WildGuard missing on two harmful responses (Qwen2.5-32B original; OBLITERATUS
\texttt{advanced}) and one benign response (Qwen2.5-7B SOM).
\item Judge revisions: WildGuard, Qwen3Guard and the HarmBench classifier identical in all cells; one encoder-based
refusal classifier, not used in any analysis here, occurs in three revisions (12, 19 and 11 cells).
\item Heretic~1.x provenance: adapter SHA-256 equal to the manifests for all twelve adapters; rebuilt merges
bit-identical to the archived merges for the eight 14B and 32B adapters (579 tensors for 14B models, 771 for 32B
models); the archived Qwen2.5-32B code merge contains only 627 of 771 tensors (all identical to the rebuild), so the
complete rebuild was evaluated. No tensor hashes of the archived 7B merges had been recorded, so the four 7B merges
were not compared tensor by tensor; for the Qwen2.5-7B general variant a separate CPU check found the archived merge
equal to the base model plus the archived adapter (all 20 edited tensors bit-identical to an in-memory merge, 40 of
40 sampled unedited tensors identical to the base).
\item Heretic~1.x settings, read from the settings record of all 12 Qwen2.5 studies: 32 keyword markers (no
``disclaimer''), system prompt ``You are a helpful coding assistant.'', 200 trials (60 random), seed~0, KL target
0.01, direction projection off, row normalisation none, at most 100 dev response tokens.
\item Harness revision: the Qwen2.5-32B original's capability and held-out runs used an earlier revision of the
evaluation script than its edited models, without the U+2581 output post-processing (its outputs contain no
U+2581) and with unpatched MBPP per-item samples (summary rescored), so MBPP changes of all Qwen2.5-32B cells,
including the KL $\le$ 0.1 re-export, are unpaired; all other cells of a base model share one revision
\cite{young2026harder}.
\item The companion study's Qwen2.5-7B and Qwen3.5-4B originals and models were generated on an RTX 4090 and judged
from identical copies (same response checksum) on the RTX PRO 6000; the new Qwen2.5-7B, Qwen2.5-Coder-7B and
Qwen3.5-4B cells were generated and judged on the RTX 4090; the 14B and 32B cells were generated and judged on the
RTX PRO 6000.
\item Teacher-forced divergence used the same reference responses and token counts for every edited model and its
original.
\item More than ten of the 972 harmful responses were flagged degenerate in three cells: OBLITERATUS
\texttt{spectral\_cascade} (56), OBLITERATUS \texttt{advanced} (30) and 1.x general on Qwen2.5-Coder-7B (11); every
other cell had at most 9.
\end{itemize}

\section{Selection-Cap Re-Export: Provenance and Validation}
\label{app:reexport}

% src (reexport_1x/VERSION.md evidence 1-6 and "Settings that are not in the journal"; results_rtx/*/PROVENANCE.json:
%      environment, datasets, base_refusals_journal/recounted, validation, validation_criteria, merge_formula_selfcheck;
%      results_rtx/reexport_1x.log l.237-262; reexport_1x/run_on_rtx.sh header; validation_7b/SUMMARY.json)
\paragraph{Version.} The 1.x studies store their settings as the serialised settings object of the Heretic release
that ran them. Its 27 keys, in their stored order, occur only in the settings class of commits from \texttt{6757ada}
(13 June 2026) to the commit before \texttt{0146b27} (27 June 2026), a range that includes the 1.4.0 release.
Within that range, commit \texttt{00185db} (18 June) lowered the search range of the MLP ablation weight from
[0.8, 1.5] to [$-$0.25, 1.5]; the journals, written from 19 to 21 June, record [0.8, 1.5], so the studies ran from
the 1.4.0 release (tag \texttt{v1.4.0}, commit \texttt{6ea3b8d}; the PyPI wheel is byte-identical to the tag) rather
than from a development checkout. The stored trial attributes and objective values, the journal format and the PEFT
version recorded in the adapters (0.19.1, the version locked by 1.4.0) agree with that release. Heretic~2.0.0
(\texttt{0db5c31}) carries the lowered range: the journals of the companion study's DiM searches record
[$-$0.25, 1.5] for the MLP maximum weight, clamped at zero, and every other search range equal to that of 1.4.0
(Table~\ref{tab:h1setup}).

\paragraph{Procedure and environment.} The re-export ran heretic-llm 1.4.0 in Python 3.12.16 with transformers
5.6.2, PEFT 0.19.1, accelerate 1.13.0, Optuna 4.8.0 and datasets 4.8.4 (the 1.4.0 lock) and PyTorch 2.14.1 with
CUDA 13.0, on an NVIDIA RTX PRO 6000 Blackwell Server Edition with TF32 matrix multiplication disabled. The direction
and evaluation prompts were the four dataset splits named in the journal (\texttt{harmful\_behaviors} and
\texttt{harmless\_alpaca}, train[:400] and test[:100]); the SHA-256 of each prompt list is recorded. Settings that
1.4.0 does not store but that change the computation were taken from the archived configuration files: the archived
runs computed residuals for all prompts and averaged them in single precision, which the re-export reproduces. The
recount of base-model dev refusals gave 97 of 100, against 98 in the journal; the 1.x objective normalises the
refusal count by this number.

\paragraph{Acceptance criteria.} The re-export of a trial was accepted as numerically equivalent if the edited
modules matched, the minimum direction cosine between corresponding adapter modules was at least 0.9995, the largest
relative adapter difference at most 0.02, the dev refusals within 5 of the archived adapter re-evaluated in the same
run and within 8 of the journal, and the dev KL within the larger of 0.002 and 15\% (archived adapter) or 30\%
(journal) of the reference value. For trial~\#84 (the released trial), the 62
edited modules matched, the minimum direction cosine was 0.99999, the largest relative differences of the two LoRA
factors were 0.0072 and 0.0050, the merged model had 709 of 771 tensors bit-identical with a maximum absolute
difference of 0.0078 and a maximum relative Frobenius difference of 0.07\% of the weight, and the dev results were
90 refusals at KL 0.0055 (re-export), 90 at 0.0054 (archived adapter, same run) and 89 at 0.0059 (journal). For
trial~\#96 no archived weights exist; it edited 80 modules, its merge passed a self-check of the merge formula on 8
tensors, and it gave 51 dev refusals at KL 0.0637 against 50 at 0.0642 in the journal. Both trials passed. Before
the GPU run, the same script had been checked on CPU against the released Qwen2.5-7B general trial (\#193): 20 of 20
edited modules, minimum direction cosine 0.99999, 319 of 339 merged tensors bit-identical (the 20 others exactly the
edited ones, maximum absolute difference 0.0020), and 65 dev refusals against 64 for the archived adapter on the same
CPU and 67 in the journal.

\paragraph{Code study.} The 2.0 rule would select trial~\#134 of the Qwen2.5-32B 1.x code search (2 of 100 dev
refusals on its malicious-code dev set at KL 0.072), with the released trial~\#179 (30 at 0.009) as its validation
reference. Its four prompt files were generated by the archived study's data-preparation script and were never
copied to the analysis server; without them the script stops at its input check, so the code study was skipped in
the run (\texttt{SKIP\_CODE=1} in the run log) and no code re-export exists.

\paragraph{Records.} \texttt{reexport\_1x/VERSION.md} (version evidence), \texttt{export\_1x\_trial.py} (re-export
script), \texttt{results\_rtx/} (provenance records of both trials, adapter of trial~\#96, tensor hashes, the run log
and the installed package list) and \texttt{validation\_7b/} (CPU check).

\section{External Tools: Provenance}
\label{app:external}

% src (/data/heretic_2/models/external/*/PROVENANCE.json; tools_ext/README.md (layout, preset mapping, overlap table,
%      runtime patches, Qwen3.5 compatibility); 33_.../runs/local-4090/logs/external_small_prompt_overlap.json;
%      tool_artifacts/prompts_used.json of the three jwest33 runs: identical lists (sha256 of JSON 9ed92df66b20 /
%      6ca52f0c8bb6); refusal_sets.json: jbb_harmful_20 and jbb_harmful_25 have advbench_overlap "exact")
\paragraph{Software.} Each tool ran in its own virtual environment with PyTorch 2.14.1 (CUDA 13.0) and transformers
5.18.0, the versions of the Heretic environment of the companion study, on the RTX 4090. OBLITERATUS was fetched at
tag \texttt{v0.1.4} (commit \texttt{461ff086}) and jwest33's abliterator at commit \texttt{6ca3356e}; both
checkouts had no local modifications. A shared finaliser converted each output to the Hugging Face layout of the
Heretic exports (bf16 shards, the base model's configuration, tokenizer and chat template; the tool's own
configuration differences, directions and metadata are kept alongside) and recorded the SHA-256 of every shard. The
provenance record of each cell stores the tool, commit, preset and its full resolved configuration, runtime patches,
command, environment, seed, base revision, prompt sets with hashes, the tool's own metrics, wall time and GPU.

\paragraph{OBLITERATUS.} All three presets ran through \texttt{AbliterationPipeline} with the preset name, seed~1
and bf16, which is what the command-line handler passes except for its float16 default; the \texttt{som} preset is
not reachable from the command line and fails in the pipeline because its direction module is not shipped. The
presets edited ten (\texttt{basic}), seven (\texttt{advanced}) and 16 (\texttt{spectral\_cascade}) layers chosen by
the tool's layer-selection heuristic and also projected the output head. The tool's own quality check, on a
verification sample of 30 prompts from its built-in harmful list, reported refusal rates of 83\%, 40\% and 3\%,
perplexity 3.98, 4.69 and 7.41 (original 3.11), coherence 0.7, 1.0 and 0.9 and no degenerate responses; it reported
a sequence-level KL divergence of 1.12, 0.51 and 1.75 nats and a first-token KL divergence of 0.081, 3.40 and 8.82.
None of these presets optimises KL, so the tool's KL budget of 0.5 was reported but not enforced. On Qwen3.5-4B the
pipeline stops at its baseline validation, which accepts hybrid Qwen surgery only for the 27-billion-parameter
topology it was validated on, before any weight is modified. An opt-in patch that generalises that check exists in
the wrapper but was not used, because results with it would lie outside the tool's validated scope.

\paragraph{jwest33 abliterator.} The batch flags were \texttt{--dtype bfloat16 --device cuda --winsorize
--null-space}; all other values were the tool's defaults (directions at the last prompt token, normalised and
norm-preserving, winsorisation at the 99.5th percentile, null-space rank ratio 0.95, refusal filter threshold
$-$7.0 on the log-likelihood of refusal over 50 tokens, hybrid strategy ``auto'' with weight 1.0 on full-attention
and 0.4 on linear-attention layers). The prompt pools from commit \texttt{a2a7940} hold 5{,}000 harmful, 5{,}000
harmless and 287 preservation prompts. With seed~1 the three runs sampled the same 30 harmful and 30 harmless
prompts. The harmful sample contains one held-out JailbreakBench prompt (also an AdvBench prompt), four AdvBench
prompts in total and two prompts that match the dev set (one exactly, one approximately); none of the 60 prompts is
a held-out benign prompt. Overlap was checked with the rule used to build the held-out set (exact after normalisation,
or token Jaccard similarity of at least 0.6 and sequence ratio of at least 0.8). The tool detected Qwen3.5-4B as
hybrid, with full attention in layers 3, 7, \ldots, 31, and Qwen2.5-7B as non-hybrid. The sensitivity patch sets this
detected information before direction extraction, as the tool's library entry point does; a second opt-in patch,
which corrects the projection side of square residual writers, was not used.

\section{Seed-Noise Threshold: Derivation}
\label{app:seednoise}

% src (2026-11-08): analysis/revision_20261108/seednoise_pairs.csv and seednoise_derivation.csv, recomputed by
% revision_analyses.py from paper-33 results/analysis/seednoise_seed_pairs.csv and review_sel_seed_runs.csv (batch
% sizes from each run's output.log); the values equal paper-33 seednoise_sigma.csv (sigma 9.1542 [6.3962, 16.0650],
% threshold 25.8920 [18.0912, 45.4387]) and review_sel_seed_sigma.csv (7 pairs: sigma 7.1965 [4.7581, 14.6468],
% threshold 20.3548 [13.4581, 41.4275]). Cell bootstrap of sigma here: [5.81, 12.11] (20,000 resamples, seed 20261108).
The threshold used for search-level verdicts (\S\ref{sec:stats}) is estimated from the ten cells of the companion study
\cite{young2026harder} that were optimised and evaluated twice with different search seeds: Heretic~2.0 DiM, MPOA and
ARA and the fork plug-ins SOM and J-space, each on Qwen2.5-7B and Qwen3.5-4B, with 100 TPE trials, the selection rule of
\S\ref{sec:models} and the held-out harness of this paper (Table~\ref{tab:seedpairs}). The seed-1 cells are the
Qwen2.5-7B and Qwen3.5-4B Heretic~2.0 and fork cells evaluated here. For each pair, $d_i$ is the seed-2 minus seed-1
WildGuard residual refusal on the same 972 harmful prompts, so prompt sampling is held fixed and $d_i$ reflects the
search. Assuming a mean difference of zero (observed mean $+0.13$~pp), the run-to-run standard deviation of one cell is
\[
\hat\sigma=\sqrt{\frac{1}{2n}\sum_{i=1}^{n} d_i^2}=9.2~\text{pp}, \qquad n=10,
\]
with a 95\% interval from $\sum d_i^2/(2\sigma^2)\sim\chi^2_n$ of [6.4, 16.1]~pp (a bootstrap over the ten cells gives
[5.8, 12.1]). The standard deviation of a difference between two independently optimised cells is
$\sqrt{2}\,\hat\sigma=12.9$~pp [9.0, 22.7], and the threshold of two such standard deviations is 25.9~pp [18.1, 45.4],
the interval propagated from that of $\hat\sigma$. Median $|d_i|$ was 9.4~pp and the largest 25.2~pp (SOM on
Qwen3.5-4B).

Three seed-2 runs on Qwen3.5-4B (DiM, SOM and ARA) fell back to a batch size of 64 after running out of memory, while
their seed-1 runs and all other runs used 128, so those pairs differ in batch size as well as in seed. Without them
($n=7$), $\hat\sigma=7.2$~pp [4.8, 14.6] and the threshold is 20.4~pp [13.5, 41.4]. The search-level rule uses the
upper limit of the ten-pair interval (45.4~pp), which none of the version contrasts reaches; under the point threshold
of either estimate the H1 contrast on Qwen2.5-32B (25.9~pp) and the rule-matched contrast with 2.0 DiM (35.8~pp) would
count as exceeding search noise and the H1 contrast on Qwen2.5-7B (19.0~pp) would not.

Transfer to the contrasts of this paper is an assumption. The replicates come from 100-trial Heretic~2.0 and fork
searches on models of 4 and 7 billion parameters; the 1.x arm used 200-trial searches with another objective, and the
reversal concerns Qwen2.5-32B, for which no replicate exists. A larger trial budget may reduce the variation of the
selected trial, and a search that stalls, as the Qwen2.5-32B DiM search did at 89 dev refusals, may increase it.
Neither direction was tested; a second seed per arm on Qwen2.5-32B would replace the assumption with an estimate.

\begin{table}[!htbp]
\centering\small
\caption{Seed replicates of the companion study \cite{young2026harder} used for the seed-noise threshold. Residual
refusal: WildGuard, 972 held-out harmful prompts (\%). $d$ = seed 2 minus seed 1 (pp). Batch: evaluation batch size
of the search (seed 1 / seed 2). SOM and J-space are fork plug-ins.}
\label{tab:seedpairs}
\setlength{\tabcolsep}{5pt}
\begin{tabular}{@{}llrrrc@{}}
\toprule
Model & Configuration & Seed 1 & Seed 2 & $d$ & Batch \\
\midrule
Qwen2.5-7B & 2.0 DiM & 20.6 & 4.4 & $-$16.2 & 128 / 128 \\
 & 2.0 MPOA & 13.8 & 3.6 & $-$10.2 & 128 / 128 \\
 & SOM & 12.2 & 17.5 & +5.2 & 128 / 128 \\
 & J-space & 52.6 & 61.2 & +8.6 & 128 / 128 \\
 & 2.0 ARA & 4.7 & 2.2 & $-$2.6 & 128 / 128 \\
Qwen3.5-4B & 2.0 DiM & 54.6 & 48.8 & $-$5.9 & 128 / 64 \\
 & 2.0 MPOA & 18.0 & 33.7 & +15.7 & 128 / 128 \\
 & SOM & 21.9 & 47.1 & +25.2 & 128 / 64 \\
 & J-space & 73.8 & 71.8 & $-$2.0 & 128 / 128 \\
 & 2.0 ARA & 18.9 & 2.2 & $-$16.8 & 128 / 64 \\
\bottomrule
\end{tabular}
\end{table}

\section{Relation to Prior Work}
\label{app:versions}

An earlier comparison by the same author \cite{young2025comparative} reported tool compatibility on sixteen 7--14B
models and keyword-based refusal and benchmark changes for Heretic, llm-abliteration (labelled DECCP there), ErisForge and FailSpy.
The present study is a separate paper, not a revision of that one: it uses mostly different base models (only
Qwen2.5-7B-Instruct is shared, and its cells here come from new searches, archived 1.x runs from June 2026 and new
external-tool runs, not from the earlier comparison), held-out prompts, model-based judges and paired tests, and none
of the earlier numbers is reused. All results here come from the held-out harness shared with the companion study
\cite{young2026harder}, which reports 19 of the cells evaluated here (13 in the main design and the six Qwen3.5-4B
context cells) for a different question, whether newer model generations are harder to abliterate with one tool. The
claims divide by result family. First reported in the companion study, and reused here as comparators: the Heretic~2.0
and fork-plug-in models on Qwen2.5-7B, Qwen2.5-32B and Qwen3.5-4B, the originals of those models, the two community
releases, their judge labels, the harness specification (including the degeneracy heuristic of
\S\ref{sec:measures}) and the seed replicates. First reported here: the twelve Heretic~1.x models and the four
coder and 14B originals, the selection-cap re-export, the external-tool cells, the tests of H1--H3, the dev-set
comparison at matched dev KL, the search ranges, costs and direction cosines, and the derivation of the seed-noise
threshold in Appendix~\ref{app:seednoise}. No claim of the companion study about model generations is made here.
